\section{Weighted kernel identities and large-sieve transfer}
\label{sec:weighted-kernel-large-sieve}

This section develops the coefficient-uniform part of the argument.  All
coefficients may be complex, and all interval estimates are uniform in the
locations of the intervals.  We first retain the exact finite Fourier
identities, then use two different kinds of orthogonality.  Orthogonality at
one exact conductor controls complete shift means, while the classical
multiplicative large sieve, summed over dyadic conductor blocks, controls a
fixed shift in the range stated below.

\subsection{The weighted centered form}

Let $w\geq 5$, and put
\begin{equation}\label{eq:weighted-Q-window}
 \cP(w)=\{p:w<p\leq w^2\},
 \qquad
 Q=\prod_{p\in\cP(w)}p,
 \qquad
 G=(\Z/Q\Z)^\times,
 \qquad
 g=|G|=\varphi(Q).
\end{equation}
Write $r=\omega(Q)$ and
\begin{equation}\label{eq:weighted-local-products}
 \delta_Q=\frac{g}{Q},
 \qquad
 \kappa=\prod_{p\mid Q}\frac{p-2}{p-1},
 \qquad
 \Xi=\prod_{p\mid Q}\frac{p}{p-2}.
\end{equation}
For an integer $h$ satisfying $(h,Q)=1$ and $a,b\in G$, define
\begin{align}
 K_h(a,b)
  &=\prod_{p\mid Q}\one_{ab+h\not\equiv0\pmod p},
  \label{eq:weighted-binary-kernel}\\
 K_h^\circ(a,b)
  &=\kappa^{-1}K_h(a,b)-1.
  \label{eq:weighted-centered-kernel}
\end{align}
The factor $\kappa$ is the exact row and column mean of $K_h$ on $G^2$.

By the Chinese remainder theorem, every $\chi\in\whG$ has local components
$\chi_p\in\widehat{\F_p^\times}$.  Set
\begin{equation}\label{eq:weighted-character-conductor}
 q_\chi=\prod_{\substack{p\mid Q\\\chi_p\ne\one}}p,
 \qquad
 d(\chi)=d(q_\chi)
 =\prod_{p\mid q_\chi}\frac1{p-2}.
\end{equation}
We extend characters by zero off the units.  If $\chi\ne\one$, then it is
induced by a primitive character of the squarefree conductor $q_\chi$.
Local character orthogonality gives
\begin{equation}\label{eq:weighted-exact-kernel-expansion}
 K_h^\circ(a,b)
 =\sum_{\substack{\chi\in\whG\\\chi\ne\one}}
 c_h(\chi)\chi(a)\chi(b),
 \qquad
 c_h(\chi)
 =(-1)^{\omega(q_\chi)}d(\chi)\overline{\chi(-h)}.
\end{equation}
In particular,
\begin{equation}\label{eq:weighted-exact-conductor-count}
 \#\{\chi\in\whG:q_\chi=q\}
 =\prod_{p\mid q}(p-2)=d(q)^{-1}
 \qquad (1<q\mid Q).
\end{equation}

Let $I,J\subset\Z$ be arbitrary integer intervals of respective lengths
$1\leq M,N<Q$.  Let $(a_m)_{m\in I}$ and $(b_n)_{n\in J}$ be arbitrary
complex coefficients.  We use the effective unit-support norms
\begin{align}
 L_a&=\sum_{\substack{m\in I\\(m,Q)=1}}|a_m|,
 &E_a&=\sum_{\substack{m\in I\\(m,Q)=1}}|a_m|^2,
 \label{eq:weighted-a-norms}\\
 L_b&=\sum_{\substack{n\in J\\(n,Q)=1}}|b_n|,
 &E_b&=\sum_{\substack{n\in J\\(n,Q)=1}}|b_n|^2,
 \label{eq:weighted-b-norms}
\end{align}
and the transforms
\begin{align}
 A_a(\chi)&=\sum_{\substack{m\in I\\(m,Q)=1}}a_m\chi(m),
 &A_b(\chi)&=\sum_{\substack{n\in J\\(n,Q)=1}}b_n\chi(n).
 \label{eq:weighted-character-transforms}
\end{align}
The weighted centered box form is
\begin{equation}\label{eq:weighted-box-form}
 \cB_h(a,b)
 =\sum_{\substack{m\in I,\ n\in J\\(mn,Q)=1}}
 a_mb_nK_h^\circ(m,n).
\end{equation}

\begin{proposition}[Weighted character expansion and shift Parseval]
\label{prop:weighted-shift-parseval}
For every $(h,Q)=1$,
\begin{equation}\label{eq:weighted-box-character-expansion}
 \cB_h(a,b)
 =\sum_{\substack{\chi\in\whG\\\chi\ne\one}}
 c_h(\chi)A_a(\chi)A_b(\chi).
\end{equation}
Moreover, for the fully active shifts $h=2\ell$,
\begin{equation}\label{eq:weighted-complete-shift-parseval}
 \frac1g\sum_{\ell\in G}|\cB_{2\ell}(a,b)|^2
 =\sum_{\substack{\chi\in\whG\\\chi\ne\one}}
 d(\chi)^2|A_a(\chi)|^2|A_b(\chi)|^2.
\end{equation}
Both identities are exact.
\end{proposition}

\begin{proof}
For $p\nmid h$ and $t\in\F_p^\times$, orthogonality on
$\F_p^\times$ gives
\[
 \frac{p-1}{p-2}\one_{t\ne-h}
 =1-\frac1{p-2}
 \sum_{\substack{\chi_p\in\widehat{\F_p^\times}\\
                   \chi_p\ne\one}}
 \overline{\chi_p(-h)}\chi_p(t).
\]
Multiplying over $p\mid Q$ and deleting the all-principal term proves
\eqref{eq:weighted-exact-kernel-expansion}.  Substitution into
\eqref{eq:weighted-box-form}, followed by finite rearrangement, proves
\eqref{eq:weighted-box-character-expansion}.

For $h=2\ell$,
\[
 c_{2\ell}(\chi)
 =(-1)^{\omega(q_\chi)}d(\chi)
   \overline{\chi(-2)}\,\overline{\chi(\ell)}.
\]
After expanding the square in
\eqref{eq:weighted-box-character-expansion}, the normalized average over
$\ell\in G$ of $\overline{\chi(\ell)}\psi(\ell)$ is
$\one_{\chi=\psi}$.  This proves
\eqref{eq:weighted-complete-shift-parseval}; complex coefficients cause no
change because the second copy of the box form is conjugated.
\end{proof}

\subsection{Deterministic completion of a shift interval}

Let $\cL$ be an arbitrary interval of $H$ consecutive integers and set
\begin{equation}\label{eq:weighted-shift-interval-count}
 H_Q=\#\{\ell\in\cL:(\ell,Q)=1\}.
\end{equation}
Inclusion--exclusion gives
\begin{equation}\label{eq:weighted-HQ-discrepancy}
 |H_Q-H\delta_Q|\leq2^r.
\end{equation}
Define the explicit completion cost
\begin{equation}\label{eq:weighted-completion-cost}
 \mathfrak C(Q)
 =\kappa^{-2}4^r+2\kappa^{-1}3^r+2^r.
\end{equation}

\begin{proposition}[Weighted deterministic shift completion]
\label{prop:weighted-shift-completion}
For arbitrary complex coefficients as above,
\begin{align}
 \left|
 \sum_{\substack{\ell\in\cL\\(\ell,Q)=1}}
 |\cB_{2\ell}(a,b)|^2
 -\frac HQ\sum_{\ell\in G}|\cB_{2\ell}(a,b)|^2
 \right|
 \leq \mathfrak C(Q)L_a^2L_b^2.
 \label{eq:weighted-shift-completion}
\end{align}
If $H\delta_Q\geq2^{r+1}$, then
\begin{align}
 \frac1{H_Q}
 \sum_{\substack{\ell\in\cL\\(\ell,Q)=1}}
 |\cB_{2\ell}(a,b)|^2
 \ll
 \frac1g\sum_{\ell\in G}|\cB_{2\ell}(a,b)|^2
 +\frac{\mathfrak C(Q)}{H\delta_Q}L_a^2L_b^2.
 \label{eq:weighted-short-shift-transfer}
\end{align}
The implied constant is absolute.
\end{proposition}

\begin{proof}
Put
\[
 S_\ell
 =\sum_{\substack{m\in I,\ n\in J\\(mn,Q)=1}}
 a_mb_nK_{2\ell}(m,n),
 \qquad
 T=\left(\sum_{\substack{m\in I\\(m,Q)=1}}a_m\right)
   \left(\sum_{\substack{n\in J\\(n,Q)=1}}b_n\right).
\]
Then $\cB_{2\ell}=\kappa^{-1}S_\ell-T$.  In the expansion of
$|S_\ell|^2$, the unit condition on $\ell$ and the two kernel conditions
exclude at most three residue classes modulo each $p\mid Q$.  Expanding
these exclusions by inclusion--exclusion produces at most $4^r$ terms.
Each term is a single residue class modulo some divisor of $Q$, and its
count in an interval differs from $H/Q$ times its count in a complete
period modulo $Q$ by at most one.
The corresponding discrepancies for $S_\ell\overline T$ and $|T|^2$ are
at most $3^r$ and $2^r$, respectively.  The total absolute coefficient
masses of all four terms are at most $L_a^2L_b^2$.  The coefficients of
the four terms in
\[
 |\kappa^{-1}S_\ell-T|^2
 =\kappa^{-2}|S_\ell|^2
  -\kappa^{-1}S_\ell\overline T
  -\kappa^{-1}\overline{S_\ell}T+|T|^2
\]
therefore give exactly the right side of
\eqref{eq:weighted-shift-completion}.

Under the stated lower bound for $H\delta_Q$,
\eqref{eq:weighted-HQ-discrepancy} implies
$H_Q\geq H\delta_Q/2$.  Divide
\eqref{eq:weighted-shift-completion} by $H_Q$ and use
$(H/Q)g=H\delta_Q$ to obtain
\eqref{eq:weighted-short-shift-transfer}.
\end{proof}

\subsection{Exact-conductor orthogonality}

The next lemma is a fixed-modulus orthogonality estimate, not a new large
sieve inequality.  Its advantage is that its modulus is the exact conductor
$q$, rather than the ambient product $Q$.

\begin{lemma}[Second moment at one exact conductor]
\label{lem:weighted-exact-conductor-orthogonality}
For every divisor $q>1$ of $Q$,
\begin{align}
 \sum_{\substack{\chi\in\whG\\q_\chi=q}}|A_a(\chi)|^2
 &\leq(M+q)E_a,
 \label{eq:weighted-exact-q-second-moment-a}\\
 \sum_{\substack{\chi\in\whG\\q_\chi=q}}|A_b(\chi)|^2
 &\leq(N+q)E_b.
 \label{eq:weighted-exact-q-second-moment-b}
\end{align}
The estimates are uniform in the interval locations.
\end{lemma}

\begin{proof}
Characters with $q_\chi=q$ are induced by primitive characters modulo $q$,
so they form a subset of all Dirichlet characters modulo $q$.  Put
$u_m=a_m\one_{(m,Q)=1}$.  On the support of $u_m$, the inflated character
$\chi$ equals its inducing character modulo $q$.  Orthogonality modulo $q$
therefore gives
\[
 \sum_{\chi\bmod q}
 \left|\sum_{m\in I}u_m\chi(m)\right|^2
 =\varphi(q)
 \sum_{\substack{c\bmod q\\(c,q)=1}}
 \left|\sum_{\substack{m\in I\\m\equiv c\pmod q}}u_m\right|^2.
\]
Every residue class contains at most $M/q+1$ members of $I$.  Cauchy's
inequality inside each class, followed by $\varphi(q)\leq q$, bounds the
last display by $(M+q)E_a$.  Restricting the nonnegative sum to the exact
conductor characters proves
\eqref{eq:weighted-exact-q-second-moment-a}.  The proof of
\eqref{eq:weighted-exact-q-second-moment-b} is identical.
\end{proof}

We record the conductor masses which occur after applying the preceding
lemma.  The final estimate, which removes the one-prime supports, is useful
for prime-supported coefficients later in the paper.

\begin{lemma}[Conductor masses in the quadratic window]
\label{lem:weighted-conductor-masses}
As $w\to\infty$,
\begin{align}
 \sum_{\substack{q\mid Q\\q>1}}q\,d(q)^2&=O(1),
 \label{eq:weighted-q-d2-mass}\\
 \sum_{\substack{q\mid Q\\q>1}}d(q)^2
 &\ll\frac1{w\log w},
 \label{eq:weighted-d2-mass}\\
 \sum_{\substack{q\mid Q\\\omega(q)\geq2}}d(q)^2
 &\ll\frac1{w^2(\log w)^2}.
 \label{eq:weighted-d2-two-prime-tail}
\end{align}
Moreover, $\delta_Q=1/2+o(1)$, $\kappa=1/2+o(1)$, and
$\Xi=4+o(1)$.
\end{lemma}

\begin{proof}
Euler products give
\begin{align*}
 \sum_{q\mid Q}q\,d(q)^2
 &=\prod_{p\mid Q}\left(1+\frac{p}{(p-2)^2}\right),\\
 \sum_{q\mid Q}d(q)^2
 &=\prod_{p\mid Q}\left(1+\frac1{(p-2)^2}\right).
\end{align*}
Mertens' theorem on $w<p\leq w^2$ shows that
$\sum_{p\mid Q}p/(p-2)^2=O(1)$, proving
\eqref{eq:weighted-q-d2-mass}.  Also
\[
 s_w:=\sum_{p\mid Q}\frac1{(p-2)^2}
 \ll\frac1{w\log w}.
\]
Thus the nonempty part of the second Euler product is
$O(s_w)$, which proves \eqref{eq:weighted-d2-mass}.  The portion containing
at least two local factors is at most
$e^{s_w}-1-s_w=O(s_w^2)$, proving
\eqref{eq:weighted-d2-two-prime-tail}.  The remaining local-product
asymptotics follow from the prime number theorem and Mertens' theorem.
\end{proof}

\begin{theorem}[Coefficient-uniform complete-shift mean square]
\label{thm:weighted-complete-shift}
For arbitrary complex coefficients,
\begin{align}
 \frac1g\sum_{\ell\in G}|\cB_{2\ell}(a,b)|^2
 &\leq
 L_b^2E_a
 \sum_{\substack{q\mid Q\\q>1}}d(q)^2(M+q),
 \label{eq:weighted-complete-one-sided-a}\\
 \frac1g\sum_{\ell\in G}|\cB_{2\ell}(a,b)|^2
 &\leq
 L_a^2E_b
 \sum_{\substack{q\mid Q\\q>1}}d(q)^2(N+q).
 \label{eq:weighted-complete-one-sided-b}
\end{align}
In particular, if $|a_m|,|b_n|\leq1$, then
\begin{equation}\label{eq:weighted-complete-bounded-coefficients}
 \frac1{M^2N^2}
 \frac1g\sum_{\ell\in G}|\cB_{2\ell}(a,b)|^2
 \ll
 \frac1{\max\{M,N\}}+\frac1{w\log w}.
\end{equation}
The implied constant is absolute for all sufficiently large $w$.
\end{theorem}

\begin{proof}
In \eqref{eq:weighted-complete-shift-parseval}, group the characters by
their exact conductor and use $|A_b(\chi)|\leq L_b$.
\Cref{lem:weighted-exact-conductor-orthogonality} gives
\[
 \frac1g\sum_{\ell\in G}|\cB_{2\ell}(a,b)|^2
 \leq L_b^2E_a
 \sum_{\substack{q\mid Q\\q>1}}d(q)^2(M+q),
\]
which is \eqref{eq:weighted-complete-one-sided-a}.  Interchanging the two
variables proves \eqref{eq:weighted-complete-one-sided-b}.

For coefficients bounded by one, $L_a\leq M$, $E_a\leq M$,
$L_b\leq N$, and $E_b\leq N$.  The first inequality, divided by
$M^2N^2$, is at most
\[
 \frac1M\sum_{\substack{q\mid Q\\q>1}}q\,d(q)^2
 +\sum_{\substack{q\mid Q\\q>1}}d(q)^2.
\]
The second inequality gives the same expression with $N$ in place of $M$.
Take the smaller of the two bounds and apply
\cref{lem:weighted-conductor-masses}.
\end{proof}

For the quadratic window, the prime number theorem gives
\begin{equation}\label{eq:weighted-logQ-r-asymptotics}
 \log Q=(1+o(1))w^2,
 \qquad
 r=(1+o(1))\frac{w^2}{2\log w}.
\end{equation}
Consequently, for
\begin{equation}\label{eq:weighted-short-H-choice}
 H=\left\lfloor Q^{1/\log w}\right\rfloor,
\end{equation}
one has
\begin{equation}\label{eq:weighted-completion-decay}
 \frac{\mathfrak C(Q)}{H\delta_Q}
 =\exp\left(
  -\bigl(1-\log2+o(1)\bigr)\frac{w^2}{\log w}
 \right)=o(1).
\end{equation}

\begin{corollary}[Coefficient-uniform short-shift mean square]
\label{cor:weighted-short-shift}
Let $\cL$ be any interval of $H$ consecutive integers, with $H$ as in
\eqref{eq:weighted-short-H-choice}.  If $|a_m|,|b_n|\leq1$, then
\begin{align}
 \frac1{M^2N^2}
 \frac1{H_Q}
 \sum_{\substack{\ell\in\cL\\(\ell,Q)=1}}
 |\cB_{2\ell}(a,b)|^2
 \ll
 \frac1{\max\{M,N\}}+\frac1{w\log w}
 +\frac{\mathfrak C(Q)}{H\delta_Q}.
 \label{eq:weighted-short-bounded-coefficients}
\end{align}
In particular, the normalized root mean square is $o(1)$ whenever
$\max\{M,N\}\to\infty$.
\end{corollary}

\begin{proof}
For large $w$, \eqref{eq:weighted-completion-decay} implies
$H\delta_Q\geq2^{r+1}$.  Apply
\eqref{eq:weighted-short-shift-transfer}, use $L_a\leq M$ and
$L_b\leq N$, and then invoke
\cref{thm:weighted-complete-shift}.
\end{proof}

\subsection{A dyadic large-sieve estimate at a fixed shift}

The preceding mean-square result exploits orthogonality in the shift.  At a
fixed shift we instead sum exact conductors in dyadic blocks.  We use the
classical multiplicative large sieve in the form
\begin{equation}\label{eq:weighted-classical-multiplicative-large-sieve}
 \sum_{q\leq S}\frac q{\varphi(q)}
 \sum_{\chi\bmod q}^{*}
 \left|\sum_{n\in\mathcal I}u_n\chi(n)\right|^2
 \ll (|\mathcal I|+S^2)\sum_{n\in\mathcal I}|u_n|^2,
\end{equation}
where the star denotes primitive characters.  The estimate is uniform in
the location of the integer interval $\mathcal I$; see, for example,
\citet[Theorem~7.13]{IwaniecKowalski2004}.  We apply this classical result
only to the sparse family of squarefree conductors dividing $Q$.

For $R\geq1$, let $\cB_h^{>R}(a,b)$ denote the portion of
\eqref{eq:weighted-box-character-expansion} with $q_\chi>R$.

\begin{theorem}[Dyadic multiplicative-large-sieve bound]
\label{thm:weighted-fixed-shift-large-sieve}
Let $I,J$ be arbitrary integer intervals of lengths $M,N<Q$, let the
coefficients be arbitrary complex numbers, and let $(h,Q)=1$.  Uniformly for
$w\leq R\leq\sqrt Q$,
\begin{equation}\label{eq:weighted-fixed-shift-tail-bound}
 |\cB_h^{>R}(a,b)|
 \ll
 \sqrt{E_aE_b}
 \left\{
  \frac{\sqrt{MN}}R
  +(\log Q)(\sqrt M+\sqrt N)
  +\sqrt Q
 \right\}.
\end{equation}
Since every nontrivial conductor $q_\chi$ is larger than $w$, the full form
satisfies
\begin{equation}\label{eq:weighted-fixed-shift-full-bound}
 |\cB_h(a,b)|
 \ll
 \sqrt{E_aE_b}
 \left\{
  \frac{\sqrt{MN}}w
  +(\log Q)(\sqrt M+\sqrt N)
  +\sqrt Q
 \right\}.
\end{equation}
All implied constants are absolute for sufficiently large $w$, and the
estimates are uniform in $h$ and in the interval locations.
\end{theorem}

\begin{proof}
Partition $R<q\leq\sqrt Q$ into geometric blocks
$S<q\leq\min\{2S,\sqrt Q\}$, beginning with $S=R$.  Since
\[
 d(q)=\frac1q\prod_{p\mid q}\frac p{p-2}\leq\frac{\Xi}{q},
\]
Cauchy's inequality gives
\begin{align*}
 &\left|
 \sum_{\substack{q\mid Q\\S<q\leq\min\{2S,\sqrt Q\}}}
 \sum_{\substack{\chi\in\whG\\q_\chi=q}}
 c_h(\chi)A_a(\chi)A_b(\chi)
 \right|\\
 &\qquad\leq\frac{\Xi}{S}
 \left(
  \sum_{\substack{q\mid Q\\S<q\leq\min\{2S,\sqrt Q\}}}
  \sum_{q_\chi=q}|A_a(\chi)|^2
 \right)^{1/2}
 \left(
  \sum_{\substack{q\mid Q\\S<q\leq\min\{2S,\sqrt Q\}}}
  \sum_{q_\chi=q}|A_b(\chi)|^2
 \right)^{1/2}.
\end{align*}
Every exact-conductor character in these sums is primitive modulo $q$.
Apply \eqref{eq:weighted-classical-multiplicative-large-sieve} to
$a_m\one_{(m,Q)=1}$ and $b_n\one_{(n,Q)=1}$, enlarging the conductor family
to all primitive characters of modulus at most $2S$.  The block is at most
\begin{align}
 \frac{\Xi}{S}
 \sqrt{(M+4S^2)(N+4S^2)}\sqrt{E_aE_b}
 \notag\\
 \leq
 \Xi\sqrt{E_aE_b}
 \left(
  \frac{\sqrt{MN}}S+2\sqrt M+2\sqrt N+4S
 \right).
 \label{eq:weighted-one-dyadic-block}
\end{align}
Summing \eqref{eq:weighted-one-dyadic-block} over dyadic $S$ from $R$ to
$\sqrt Q$ uses
\[
 \sum_S\frac1S\ll\frac1R,
 \qquad
 \#\{S\}\ll\log Q,
 \qquad
 \sum_S S\ll\sqrt Q.
\]

It remains to treat $q_\chi>\sqrt Q$.  Character Parseval on $G$ gives
\begin{equation}\label{eq:weighted-ambient-character-parseval}
 \sum_{\chi\in\whG}|A_a(\chi)|^2=gE_a,
 \qquad
 \sum_{\chi\in\whG}|A_b(\chi)|^2=gE_b.
\end{equation}
Here the hypothesis $M,N<Q$ ensures that no two integers in either interval
occupy the same residue class modulo $Q$.  Since
$d(\chi)\leq\Xi/q_\chi\leq\Xi/\sqrt Q$, Cauchy's inequality bounds this
remaining portion by
\[
 \frac{\Xi}{\sqrt Q}\,g\sqrt{E_aE_b}
 \leq\Xi\sqrt Q\sqrt{E_aE_b}.
\]
Combining the two conductor ranges and using $\Xi=O(1)$ proves
\eqref{eq:weighted-fixed-shift-tail-bound}.  Finally, all primes dividing
$Q$ exceed $w$, so every $q_\chi>1$ exceeds $w$.  The choice $R=w$ gives
\eqref{eq:weighted-fixed-shift-full-bound}.
\end{proof}

\begin{corollary}[A fixed-shift cancellation range]
\label{cor:weighted-fixed-shift-bounded}
If $|a_m|,|b_n|\leq1$, then, uniformly for $(h,Q)=1$,
\begin{equation}\label{eq:weighted-fixed-shift-normalized-bounded}
 \frac{|\cB_h(a,b)|}{MN}
 \ll
 \frac1w
 +(\log Q)\left(\frac1{\sqrt M}+\frac1{\sqrt N}\right)
 +\sqrt{\frac Q{MN}}.
\end{equation}
Consequently, if
\begin{equation}\label{eq:weighted-fixed-shift-range-conditions}
 \frac{\min\{M,N\}}{(\log Q)^2}\longrightarrow\infty,
 \qquad
 \frac{MN}{Q}\longrightarrow\infty,
\end{equation}
then $\cB_h(a,b)=o(MN)$.  In particular, for
$M=Q^{\lambda+o(1)}$ and $N=Q^{\nu+o(1)}$ with fixed
$\lambda,\nu>0$, this holds whenever $\lambda+\nu>1$.
\end{corollary}

\begin{proof}
Under the coefficient bound, $E_a\leq M$ and $E_b\leq N$.  Substitute
these estimates in \eqref{eq:weighted-fixed-shift-full-bound} and divide by
$MN$.
\end{proof}

\begin{remark}[Strict scope of the fixed-shift estimate]
\label{rem:weighted-not-complete-type-II}
The estimate in \cref{thm:weighted-fixed-shift-large-sieve} is a bilinear
large-sieve transfer for this particular centered CRT kernel.  It requires
the ambient-volume condition visible in
\eqref{eq:weighted-fixed-shift-range-conditions}; it is not a complete
prime-sensitive Type~II theorem.  The shift-mean results retain their shift
average, while the fixed-shift result retains its length restriction.
Furthermore, $K_h$ only tests whether a product plus $h$ avoids the primes
dividing $Q$.  None of the results in this section tests whether that target
is prime, estimates $\Lambda(n)\Lambda(n+2)$, overcomes sieve parity, or
implies a twin-prime lower bound.
\end{remark}
